% !TEX root = ../tate.tex

\section{Preliminaries}\label{Preliminaries}

Let \(\k\) be a field and let \(G\) be a finite abelian group.

\subsection{\(\k[G]\)-modules}\label{Gmodules}

First, we give our conventions for defining the \(G\)-tensor and the linear dual of (left) \(\k[G]\)-modules.\anibal{I suspect that the notion is G-Hom with linear obtained by mapping to a specific module to the unit of the G-tensor. It is probably a closed symmetric monoidal category, which is what the appedix might explain.}
We refer to \cref{Groups} for a more detailed view on these algebraic notions.

\begin{definition}
	Let \(M\) and \(N\) be two left \(\k[G]\)-modules.
	The \defn{\(G\)-tensor product} \(M \otimes_G N\) is the quotient
	\[
	M \otimes_G N \coloneqq (M \otimes N)/(g^{-1}m \otimes n \sim m \otimes gn).
	\]
\end{definition}

We will denote elements in \(M \ot_G N\) simply by \(m \ot n\) (instead of \(m \ot_G n\)).

% Replacing \(m\) by \(gm\) in the defining relation shows that \((m \otimes n) \sim (gm \otimes gn)\), and hence
% \[
% M \otimes_G N \cong (M \otimes N)_G \coloneqq \k \otimes_G (M \otimes N),
% \]
% the \defn{coinvariants} of \(M \otimes N\) with respect to the diagonal \(G\)-action
% \[
% g \cdot (m \otimes n) = gm \otimes gn.
% \]

% There is also a diagonal \(G\)-action on \(\Hom(M,N)\) given by
% \[
% (g \cdot f)(m) = g f(g^{-1}m),
% \]
% which identifies equivariant maps with \defn{invariants}:
% \[
% \Hom_G(M,N) = (\Hom(M,N))^G \coloneqq \Hom_{\k[G]}(\k,\Hom(M,N)).
% \]
% Note that, by definition, \(\Hom_G(M,N) = \Hom_{\k[G]}(M,N)\).
%
% If \(A\) is a \(\k\)-vector space with the trivial \(G\)-action, then there are natural identifications
% \[
% \Hom(M_G,A) \cong \Hom_G(M,A), \qquad \Hom_G(A,M) \cong \Hom(A,M^G),
% \]
% relating invariants and coinvariants.

\begin{definition}\label{defdual}
	Let \(M\) be a left \(\k[G]\)-modules.
	The \defn{linear dual}
	\[
	M^\dd \coloneqq \hom_{\k[G]}(M,\k[G])
	\]
	is again a left \(\k[G]\)-modules, with left \(G\)-action given by \((g \cdot f)(m) = f(m)g^{-1}\).
\end{definition}

Thus, if \(u^\dd\) is the dual of a basis element \(u\) of \(M\), we have that \(g \cdot u^\dd = (gu)^\dd\).

% Another option: defining the dual as a right-module and then showing how to construct a left action
%
% \begin{definition}
	% 	Let \(W\) be a chain complex of left \(\k[G]\)-modules.
	% 	The \defn{linear dual}
	% 	\[
	% 	W^\dd \coloneqq \hom_{\k[G]}(W,\k[G])
	% 	\]
	% 	is a chain complex of right \(\k[G]\)-modules with right \(G\)-action given by \((f \cdot g)(a) = f(a)g\).
	% \end{definition}
%
% Thus, if \(u^\dd\) is the dual of a basis element of \(W\), the cyclic group acts backwards, in the sense that \(u^\vee \cdot g = (g^{-1}u)^\dd\).
%
% Via the group anti-automorphism \(g \mapsto g^{-1}\), we can endow \(W^\dd\) with a left \(G\)-action, defined as
% \[
% g \cdot u^\dd \coloneqq u^\dd \cdot g^{-1} = (gu)^\dd.
% \]
% Note that \((g \cdot f)(a) = f(a)g^{-1}\), not \(gf(a)\).
% We will always consider these left actions for the dual chain complexes.

\subsection{Derived categories}

Let \(R\) denote either \(\k\) or \(\k[G]\).
Let \(\Ch(R)\) be the category of unbounded chain complexes of left \(R\)-modules equipped with the projective model structure, where weak equivalences are quasi-isomorphisms and fibrations are degreewise surjections.

\begin{definition}
	The \defn{derived category} \(\mathcal D(R)\) is the localization of \(\Ch(R)\) at the class \(\mathcal W\) of quasi-isomorphisms:
	\[
	\mathcal D(R) = \Ch(R)[\mathcal W^{-1}].
	\]
\end{definition}
\martin{add morphisms of the derived category here or in appendix.}

\begin{remark}\label{DFpmorph}
	Since \(\k\) is a field, the derived category \(\cD(\k)\) has the property that any chain complex \(V\) is canonically isomorphic to the complex \(H_*(V)\) with zero differential.
	Furthermore, any morphism in \(\mathcal D(\k)\) is completely determined by its induced map on homology (this is explained in detail in \cref{derived}).
\end{remark}

The suspension functor \([1]\) is induced by the shift of complexes and is an autoequivalence of \(\mathcal D(R)\).
Moreover, homotopy pushouts and homotopy pullbacks coincide in \(\mathcal D(R)\).
In particular, \(\hofib(f) \cong \hocofib(f)[-1]\) for any map \(f\) in \(\mathcal D(R)\).

We refer to sequences \(X \to Y \to \hocofib(f)\) and \(\hofib(f) \to X \to Y\) as cofibre sequences and fibre sequences, respectively.
Note that, by the above, they carry equivalent information.

The mapping cone and mapping cocone constructions give models for the homotopy cofiber and fiber, respectively.

\begin{definition}
	Let \(M=(M_*,d_*)\) and \(N=(N_*,d_*)\) in \(\cD(R)\), and let \(f \colon M \to N\) be a chain map. The \defn{mapping cone} \(\cone(f)\) is the chain complex
	\[
	\cone(f)_n \coloneqq N_n \oplus M_{n-1},
	\]
	with differential \(d((n,m)) \coloneqq (f(m)+dn, -dm)\).
	The \defn{mapping cocone} is \(\cone(f)[-1]\). 
\end{definition}


% The category \(\Ch(K)\) is pointed, with zero object the complex \(0\).
% \begin{definition}
	% 	For a map \(f \colon X \to Y\) in \(\Ch(K)\), the \defn{cofibre} \(\cofib(f)\) and \defn{fibre} \(\fib(f)\) are defined by the pushout and pullback squares
	% 	\[
	% 	\begin{tikzcd}
		% 		X \ar[r, "f"] \ar[d] & Y \ar[d] \\
		% 		0 \ar[r] & \cofib(f)
		% 	\end{tikzcd}
	% 	\qquad\text{and}\qquad
	% 	\begin{tikzcd}
		% 		\fib(f) \ar[r] \ar[d] & 0 \ar[d] \\
		% 		X \ar[r, "f"] & Y
		% 	\end{tikzcd}
	% 	\]
	% 	respectively.
	% \end{definition}
	
We now extend the definitions from \cref{Gmodules} to \(\cD(\k[G])\).

\begin{definition}	
	Let \(M = (M_*,d_*)\), where \(d_n \colon M_{n+1} \to M_n\), be a chain complex in \(\cD(\k[G])\). The linear dual is defined by setting
	\[
	(M^\vee)_n = (M_{-n})^\vee,
	\]
	with differential \((d^\vee)_n \colon (M^\vee)_n \to (M^\vee)_{n-1}\) given by
	\[
	(d^\vee)_n (f) \coloneqq (-1)^n (d_{-n})^\vee (f) = (-1)^n f \circ d_{-n},
	\]
	for \(f \in (M^\vee)_n\).
\end{definition}
Observe that \(d^\vee\) is indeed \(\k[G]\)-linear, with respect to the left action of \cref{defdual}:\martin{Is this worth writting here?}
\[
\left (g \cdot (d^\vee(f)\right ))(m) = \left (g \cdot (f \circ d)\right )(m) = f(dm)g^{-1} = (g\cdot f)(dm) = \left ((g \cdot f) \circ d\right )(m) = \left (d^\vee(g \cdot f)\right )(m).
\]
\begin{definition}
	If \(M\) and \(N\) in \(\cD(\k[G])\), the \(G\)-tensor \(M \otimes_G N\) is the chain complex in \(\cD(\k)\) defined as
	\[
	(M \otimes_G N)_n \coloneqq \bigoplus_{i+j=n} M_i \otimes_G N_j,
	\]
	with differential
	\[
	d(m \ot n) = (d m \ot n) + (-1)^{|m|} (m \ot dn).
	\]
\end{definition}
Note that it is well defined, \(d(g^{-1}m \ot n) = d(m \ot gn)\).
%\[
%d(g^{-1}m \ot n) = (dg^{-1}m) \ot n + (-1)^{|m|} (g^{-1}m \ot dn) = (g^{-1}dm) \ot n + (-1)^{|m|}(m \ot gdn) = (dm \ot gn) + (-1)^{|m|}(m \ot dgn) = d(m \ot gn).
%\]


\subsection{The Tate construction}

\begin{definition}
	The functors
	\[
	(-)_G = \k \otimes_G - \qquad\text{and}\qquad (-)^G = \Hom_G(\k,-)
	\]
	from \(\Ch(\k[G])\) to \(\Ch(\k)\) induce the \defn{derived coinvariants} and \defn{derived invariants} functors from \(\mathcal D(\k[G])\) to \(\mathcal D(\k)\).
	Explicitly, for \(M \in \mathcal D(\k[G])\) these are given by
	\[
	M_{hG} \coloneqq W \otimes_G M \qquad\text{and}\qquad M^{hG} \coloneqq \Hom_G(W,M),
	\]
	where \(W \xrightarrow{\sim} \k\) is a cofibrant resolution.
\end{definition}

The homology groups of \(M_{hG}\) and \(M^{hG}\) recover group homology \(H_*(G;M)\) and group cohomology \(H^*(G;M)\), respectively.

\begin{definition}
	The \defn{norm map} is the natural transformation
	\[
	\Nm \colon (-)_{hG} \to (-)^{hG}
	\]
	that derives the classical map \(M_G \to M^G\) given by the action of the norm element \(N = \sum_{g \in G} g \in \k[G]\).
\end{definition}

Concretely, since \(\Hom_{\k[G]}(\k,\k[G]) \cong \k\), the linear dual \(W^\vee = \Hom_{\k[G]}(W,\k[G])\) is itself a cofibrant resolution of \(\k\).
Thus, any chain homotopy equivalence \(W \xrightarrow{\sim} W^\vee\) induces the norm map via
\[
W \otimes_G M \longrightarrow W^\vee \otimes_G M \cong \Hom_G(W,M),
\]
where the final isomorphism holds because \(W\) consists of finitely generated projectives.

\begin{definition}
	The \defn{Tate construction} is the functor
	\[
	(-)^{tG} \defeq \hocofib(\Nm) \colon \mathcal D(\k[G]) \to \mathcal D(\k).
	\]
\end{definition}

The homology of \(M^{tG}\) recovers the \defn{Tate cohomology groups} \(\widehat H^*(G;M)\) when \(M\) is concentrated in degree \(0\).
\anibal{The word agrees makes me expect a reference.}\martin{Is recover better?}
The resulting cofibre sequence
\[
M_{hG} \xrightarrow{\Nm} M^{hG} \longrightarrow M^{tG}
\]
in \(\mathcal D(\k)\) relates group homology, group cohomology, and Tate cohomology.
We refer to it as the \defn{Tate cofibre sequence} for \(M\).
Analogously, there is a \defn{Tate fibre sequence}
\[
M^{tG}[-1] \longrightarrow M_{hG} \xrightarrow{\Nm} M^{hG}
\]
relating Tate homology, group homology and group cohomology.

\subsection{Explicit models for \(\Fp[\cyc_p]\)}

Let \(p\) be a prime and \(\cyc_p = \angles{\rho \mid \rho^p}\) be the cyclic group of order \(p\).
We work over the field \(\Fp\) with \(p\) elements and denote \(R = \Fp[\cyc_p]\).
Set \(T = \rho - 1\) and \(N = 1 + \rho + \dots + \rho^{p-1}\) in \(R\), as introduceed before.

A cofibrant resolution \(W_p \xrightarrow{\sim} \Fp\) is given by the complex
\[
W_p = \quad R\set{e_0} \xla{T} R\set{e_1} \xla{N} R\set{e_2} \xla{T} R\set{e_3} \xla{N} \dots,
\]
where each \(R\set{e_i}\) is a free \(R\)-module of rank one concentrated in homological degree \(i \geq 0\).
Since each \(R\set{e_i} \cong R\) is finitely generated and free, the linear dual is the complex
\[
W^\vee_p = \dots \xla{-N} R\set{e_3^\vee} \xla{T^\vee} R\set{e_2^\vee} \xla{-N} R\set{e_1^\vee} \xla{T^\vee} R\set{e_0^\vee},
\]
where \(T^\vee = (\rho^{-1}-1)\), and \(e_i^\dd\) has degree \(-i\).

Let \(N_{W_p} \colon W_p \xrightarrow{\sim} W^\vee_p\) be the map that is \(0\) in positive and negative degrees and sends \(e_0\) to \(N \cdot e_0^\vee\).
Its mapping cone is the complex
\[
\dots \xla{T^\vee} R\set{e_2^\vee} \xla{-N} R\set{e_1^\vee} \xla{T^\vee} R\set{e_0^\vee} \xla{N} R\set{e_0} \xla{T} R\set{e_1} \xla{N} \dots
\]
with \(e_i^\vee\) in degree \(-i\), and \(e_i\) in degree \(i+1\).
We denote by \(\widehat W_p\) the complex obtained by relabeling the generators to match their degree:
\[
\dots \xla{T^\vee} R\set{\widehat e_{-2}} \xla{-N} R\set{\widehat e_{-1}} \xla{T^\vee} R\set{\widehat e_0} \xla{N} R\set{\widehat e_1} \xla{T} R\set{\widehat e_2} \xla{N} \dots
\]

Consider the cofibre sequence
\[
W_p \xra{N} W^\vee_p \to \widehat W_p.
\]
For any chain complex \(M\) in \(\mathcal D(R)\), applying the functor \(- \ot_{\cyc_p} M\) to it gives a cofibre sequence representing the Tate cofibre sequence
\[
M_{h\cyc_p} \xra{\Nm} M^{h\cyc_p} \to M^{t\cyc_p}.
\]
In particular, \(\widehat W_p \otimes_{\cyc_p} M\) is a model for \(M^{t\cyc_p}\).

\begin{remark}
	In \cref{multiplications} we give an explicit chain level formula for a multiplication on \(\widehat W_p\), which therefore induces the cup product on Tate cohomology. % at least with \(\Fp\) coefficients.
\end{remark}

Let \(\widetilde W_p\) be the mapping cocone of the norm map.
Explicitly, \(\widetilde W_p = \widehat W_p[-1]\) is the complex
\[
\dots \xla{-N} R\set{\widetilde e_{-2}} \xla{T^\vee} R\set{\widetilde e_{-1}} \xla{N} R\set{\widetilde e_0} \xla{T} R\set{\widetilde e_1} \xla{N} R\set{\widetilde e_2} \xla{T} \dots,
\]
with \(\widetilde e_i \coloneqq \widehat e_{i+1}[-1]\) in degree \(i\).
Since it fits in the fibre sequence
\[
\widetilde W_p \to W_p \xra{N} W^\vee_p,
\]
we have that \(\widetilde W_p \otimes_{\cyc_p} M \simeq M^{t\cyc_p}[-1]\) for any \(M \in \cD(R)\).
\anibal{\(\simeq\) vs \(\cong\), worth making explicit}

\begin{remark}
	There are some important cases for which \(M^{t\cyc_p}[-1] \simeq M^{t\cyc_p}\).
	Let \(M = \Fp\), then \(H(\Fp^{t\cyc_p})\) is precisely the Tate cohomology of the cyclic group with trivial coefficients \(\widehat H^*(\cyc_p;\Fp)\), and is given by
	\[
	\widehat H^i(\cyc_p;\Fp) \cong H_i(\widehat W_p \otimes_{\cyc_p} \Fp) = H_i((\widehat W_p)_{\cyc_p}) \cong \Fp\set{\widehat e_i}.
	\]
	Similarly, we have that \(H_i(\Fp^{t\cyc_p}[-1]) = H_i(\widetilde W_p \otimes_{\cyc_p} \Fp) \cong \Fp\set{\widetilde e_i}\), which means that the isomorphism \(\widetilde e_i \mapsto \widehat e_i\) induces an equivalence \(\Fp^{t\cyc_p}[-1] \simeq \Fp^{t\cyc_p}\) in \(\cD(\Fp)\).
	\cref{fibre_cofibre} will sow that this is also the case for \(M = V^{\ot p}\) for \(V\) in \(\cD(\Fp)\). 
\end{remark}

\begin{remark}
	We can also regard \(H(\Fp^{t\cyc_p})\) as a free one-dimensional \(\Fp[u,u^{-1}]\)-module with action
	\[
	u \cdot \widehat e_i \coloneqq \widehat e_{i+1}.
	\]
\end{remark}

\subsection{Tate spectral sequences}\martin{relocate this section}

Using \(\widehat W_p\) we can prove the following lemma:

\begin{lemma}\label{spectral_seq}
	Let \(M\) in \(\mathcal D(\Fp[\cyc_p])\).
	There is a spectral sequence
	\[
	\widehat H^*(\cyc_p;H_*(M)) \Rightarrow \widehat H^*(\cyc_p;M) = H_*(M^{t\cyc_p}).
	\]
\end{lemma}

\begin{proof}
	Since \(H_*(M^{t\cyc_p}) \cong H_*(\widehat W_p \otimes_{\cyc_p} M)\), there is a filtration that gives rise to a spectral sequence with first page \(E^1 = \widehat W_p \ot_{\cyc_p}H_*(M)\).
	The second page is then \(E^2 = H_*(\widehat W_p \ot_{\cyc_p}H_*(M))\), which is isomorphic to \(\widehat H_*(\cyc_p;H_*(M))\).\qedhere
\end{proof}

\subsection{Spectra}

\begin{definition}
	The stable \(\infty\)-category of spectra.
\end{definition}

\begin{definition}
	Smash product.
\end{definition}

\begin{definition}
	Ring spectra.
\end{definition}

\subsection{The spectral tate diagonal}

Let \({\mathsf{Sp}}\) denote the stable \(\infty\)-category of spectra, and \(\wedge = \wedge_{\bS}\) its symmetric monoidal product.
In [Nikolaus-Scholze], they define the following:

\begin{definition}
	Let \(X\) be a spectrum.
	Let \(T^{\mathsf{Sp}}_p\) denote the functor \(\mathsf{Sp} \to \mathsf{Sp}\) given by
	\[
	T^{\mathsf{Sp}}_p(X) \coloneqq (X^{\wedge p})^{t\cyc_p} = \cofib(Nm \colon (X \wedge \overset{p}{\dots} \wedge X)_{h\cyc_p} \to (X \wedge \overset{p}{\dots} \wedge X)^{h\cyc_p}),
	\]
	the Tate construction of \(X^{\wedge p}\), with \(\cyc_p\)-action given by permuting the factors.
\end{definition}

They prove that it is an exact functor and give

\begin{proposition}
	For any \(F \in \Fun^{\mathrm{Ex}}(\mathsf{Sp},\mathsf{Sp})\), evaluation at the sphere spectrum \(\bS\) induces an equivalence
	\[
	\mathrm{Map}_{\Fun^{\mathrm{Ex}}(\mathsf{Sp},\mathsf{Sp})}(\id_{\mathsf{Sp}},F) \longrightarrow \Hom_{\mathsf{Sp}}(\bS,F(\bS)) = \Omega^\infty F(\bS).
	\]
\end{proposition}

Thus, the space of natural transformations \(\id_{\mathsf{Sp}} \to T_p^{\mathsf{Sp}}\) is equivalent to \(\Hom_{\mathsf{Sp}}(\bS,T_p^{\mathsf{Sp}}(\bS)) = \Omega^\infty T^{\mathsf{Sp}}_p(\bS)\).
This allows them to define

\begin{definition}
	The \defn{spectral Tate diagonal} is the natural transformation \(\Delta_p^{t,\mathsf{Sp}} \colon \id_{\mathsf{Sp}} \to T_p^{\mathsf{Sp}}\) that corresponds to the composition
	\[
	\bS \to \bS^{h\cyc_p} \to \bS^{t\cyc_p} = T_p^{\mathsf{Sp}}(\bS).
	\]
\end{definition}


\subsection{\(\sH\Fp\)-modules}\label{HFpmod}

The derived category \(\cD(\Fp)\) is equivalent to the subcategory \(\sH\Fp\text{-}\mathrm{Mod} \subset \mathsf{Sp}\) of \(\sH\Fp\)-modules, or \(\sH\Fp\)-linear spectra, via the Eilemberg-MacLane functor
\[
\sH \colon \cD(\Fp) \longrightarrow \sH\Fp\text{-}\mathrm{Mod},
\]
that sends a chain complex \(V\) in \(\cD(\Fp)\) to the spectrum \(\sH V = \bigvee_i X_i\), where each \(X_i\) is the Eilemberg-MacLane spectra with \(\pi_i X_i = H_i(V)\) and \(\pi_jX_i = 0\) if \(i \neq j\).\martin{Or is it better to say \(\sH V = \bigvee_i \Sigma^i\sH(H_i(V))\)? The thing is that \(\sH\) is used to define \(\sH\).}
On the other direction, the functor
\[
\pi_* \colon \sH\Fp\text{-}\mathrm{Mod} \longrightarrow \cD(\Fp)
\]
takes a \(\sH\Fp\)-module \(X\) to the chain complex with \(\pi_i(X)\) in degree \(i\) and zero differential.

\(\sH\) is a strong symmetric monoidal functor with respect to the tensor products \(\ot_\Fp\) of \(\cD(\Fp)\) and \(\ot_{\sH\Fp}\) of \(\sH\Fp\text{-}\mathrm{Mod}\), that is, there is an isomorphism
\[
\sH V \ot_{\sH\Fp} \sH V \xra{\cong} \sH(V \ot_\Fp V).
\]
Also, the forgetful functor \(\sH\Fp\text{-}\mathrm{Mod} \to \mathsf{Sp}\) has a lax symmetric monoidal structure given by the canonical projection
\[
\sH V \wedge_{\bS} \sH V \to \sH V \ot_{\sH\Fp} \sH V.
\]
% If \(X\) is a \(\sH\Fp\)-module, then \(T_p^{\mathsf{Sp}}(X)\) is also a \(\sH\Fp\)-module.
Their composition induces a (natural) map that Nikoulaus and Scholze denote by \(\can\)
\[
\can \colon T_p^{\mathsf{Sp}}(\sH V) = (\sH V \wedge_{\bS} \overset{p}{\cdots} \wedge_{\bS} \sH V)^{t\cyc_p} \longrightarrow \left(\sH(V \ot_\Fp \overset{p}{\cdots} \ot_\Fp V)\right)^{t\cyc_p}.
\]
Furthermore, both \(\sH\) and the forgetful functor preserve all limits and colimits, so the Tate construction in spectra is compatible to the algebraic version:
\[
\left(\sH(V \ot_\Fp \overset{p}{\cdots} \ot_\Fp V)\right)^{t\cyc_p} \cong \sH\left((V \ot_\Fp \overset{p}{\cdots} \ot_\Fp V)^{t\cyc_p}\right) = \sH(T_pV),
\]
thus, we will instead write \(\can \colon T_p^{\mathsf{Sp}}(\sH V) \to \sH(T_pV)\).

\begin{remark}
	Behavior of \(\can\) with respect to the spectral ring structure
\end{remark}

% Is the \(\cyc_p\)-action on \(H_*(M)\) trivial?
% If so we can say a bit more.

% write about the other filtration.

% \begin{lemma}
	% 	There is a spectral sequence \martin{Maybe it is better to write it without degrees and as \(H(T_2H(V);\bF_2) \Rightarrow H(T_2V;\bF_2)\).}
	% 	\[
	% 	\widehat H_p(\cyc_2;H_q(V^{\ot 2})) \Rightarrow \widehat H_{p+q}(\cyc_2;V^{\ot 2}) = H_{p+q}(T_2V;\bF_2).
	% 	\]
	% \end{lemma}
%
% \begin{proof}
	% 	Let \(W^t \coloneqq \fib(N_W \colon W \xrightarrow{\sim} W^\vee)\), the fibre of the norm map.
	% 	We have that
	% 	\[
	% 	T_2V_n \simeq \bigoplus_{p+q=n} W^t_p \ot_{\cyc_2} (V^{\ot 2})_q.
	% 	\]
	% 	Thus, there is a filtration that gives rise to a spectral sequence with first page \(E^1_{p,q} = W^t_p \ot_{\cyc_2}H_q(V^{\ot 2})\).
	% 	The second page is then \(E^2_{p,q} = H_p(W^t \ot_{\cyc_2}H_q(V^{\ot 2}))\), which is isomorphic to \(\widehat H_p(\cyc_2;H_q(V^{\ot 2}))\).
	% \end{proof}

% \subsection{Diagonals}
%
% Let \(C\) be a chain complex.
% An \defn{unstable \(p\)-cyclic diagonal} on \(C\), or \defn{unstable \(p\)-diagonal} for short, is a \(\cyc_p\)-equivariant chain map
% \[
% \mu \colon W(p) \otimes C \to C^{\otimes p},
% \]
% where the \(\cyc_p\)-action on the domain is given by multiplication on the first factor and on the codomain consists of cyclically permuting the factors.
%
% The fibre of \(N_{W(p)}\) is isomorphic to the desuspension of \(\widehat W(p)\), which we denote by
% \[
% W^{\mathrm st}(p) \coloneqq \widehat W(p)[-1] = \dots \xla{T} R\set{e_{-1}} \xla{N} R\set{e_0} \xla{T} R\set{e_1} \xla{N} R\set{e_2} \xla{T} \dots
% \]
% with \(e_i\) in degree \(i\).
% An \defn{stable \(p\)-cyclic diagonal} on \(C\), or simply a \defn{stable \(p\)-diagonal}, is a \(\cyc_p\)-equivariant chain map
% \[
% \mu \colon W^{\mathrm st}(p) \otimes C \to C^{\otimes p}.
% \]
%
% We can recover both notions of diagonals from the cofibre sequence \eqref{CofibreM} by setting \(M = \Hom(C,C^{\otimes p})\), which gives
% \[
% W(p) \otimes_R \Hom(C,C^{\otimes p}) \xra{\Nm} W^\vee(p) \otimes_R \Hom(C,C^{\otimes p}) \to \widehat W(p) \otimes_R \Hom(C,C^{\otimes p}).
% \]
% Since \(W\) consists of finitely generated projectives, we can express it in terms of homs
% \[
% \Hom_R(W^\vee(p),\Hom(C,C^{\otimes p})) \xra{\Nm} \Hom_R(W(p),\Hom(C,C^{\otimes p})) \to \Hom_R(W^{\mathrm st}(p),\Hom(C,C^{\otimes p})),
% \]
% and by hom-tensor adjunction %%% This is not clear yet
% \[
% \Hom(W^\vee(p) \otimes_R C,C^{\otimes p}) \xra{\Nm} \Hom(W(p) \otimes_R C,C^{\otimes p})) \to \Hom(W^{\mathrm st}(p) \otimes_R C,C^{\otimes p})).
% \]
% Thus, unstable and stable \(p\)-cyclic diagonals are, respectively, \(0\)-cycles in the second and third chain complexes of the cofibre sequence.